\section{Выводы}

В ходе работы решены задачи, поставленные в лабораторной работе №1:

\begin{enumerate}
  \item Написана консольная программа семейства~A для варианта~2: интерактивный
        справочник, работающий с линейным динамическим массивом через текстовое
        меню.
  \item Вся логика вынесена в библиотеку чистых функций
        \texttt{array\_ops}, использующую только \texttt{new[]} /
        \texttt{delete[]}. Запрещённые конструкции --- классы, конструкторы,
        семантика копирования и перемещения, \texttt{std::string},
        \texttt{std::vector}, \texttt{std::array}, умные указатели --- не
        используются.
  \item Реализованы сортировка выбором по возрастанию и разворот массива двумя
        указателями; на контрольном примере задания получен ожидаемый
        результат \texttt{\{3, 1, 4, 1, 5\}}
        $\rightarrow$ \texttt{\{1, 1, 3, 4, 5\}}
        $\rightarrow$ \texttt{\{5, 4, 3, 1, 1\}}.
  \item Обеспечены требования к корректности: отсутствие утечек памяти за счёт
        схемы «новая память, копирование, освобождение старой» и обнуления
        указателя, проверка \texttt{std::cin} с восстановлением состояния
        потока, штатная обработка конца ввода, корректная работа с массивом
        длины~0 и перехват \texttt{std::bad\_alloc} с сохранением исходных
        данных.
  \item Собран проект CMake (C++20) и написано 26 тестов Google Test в трёх
        наборах, подключённых к CTest, --- при требовании не менее~10 тестов.
\end{enumerate}

Практический результат: получен рабочий, собираемый и полностью покрытый
тестами пример ручного управления динамической памятью через указатели,
ссылки и out-параметры.
